% Generated draft: requires mathematical and source review.
\documentclass[11pt]{article}
\usepackage[margin=25mm]{geometry}
\usepackage{fontspec}
\setmainfont{Latin Modern Roman}
\usepackage{amsmath,amssymb,mathtools,unicode-math}
\setmathfont{Latin Modern Math}
\usepackage{enumitem,needspace,xurl,hyperref}
\setlength{\parindent}{0pt}
\setlength{\parskip}{5pt}
\setlength{\emergencystretch}{3em}
\newcommand{\N}{\mathbb{N}}
\newcommand{\Z}{\mathbb{Z}}
\newcommand{\Q}{\mathbb{Q}}
\newcommand{\R}{\mathbb{R}}
\newcommand{\C}{\mathbb{C}}
\newcommand{\F}{\mathbb{F}}
\newcommand{\A}{\mathbb{A}}
\newcommand{\PP}{\mathbb{P}}
\newcommand{\E}{\mathbb{E}}
\newcommand{\eps}{\varepsilon}
\newcommand{\dd}{\,\mathrm{d}}
\begin{document}
\section*{UnsolvedMath ID: 30001678}
\textbf{Source problem code:} catalogue entry
\subsection{Definitions and mathematical statement}
Let \(\omega=\{0,1,2,\ldots\}\) and equip \(X=\R^\omega\) with its product topology and Borel structure. A Borel equivalence relation on \(X\) is an equivalence relation \(E\subseteq X\times X\) that is Borel. A nonempty perfect set in \(\R\) is a nonempty closed set with no isolated points.

For equivalence relations \(E\) on \(X\) and \(F\) on a standard Borel space \(Y\), write \(E\leq_B F\) if there is a Borel map \(r:X\to Y\) such that
\[
\forall x,x'\in X,\qquad x\mathrel E x'\ \Longleftrightarrow\ r(x)\mathrel F r(x').
\]
For a Borel action of a Polish group \(G\) on a standard Borel space \(Y\), its orbit equivalence relation is
\[
y\mathrel{E_G^Y}y'\quad\Longleftrightarrow\quad (\exists g\in G)\;g\cdot y=y'.
\]
Here a Polish group is a topological group whose topology is separable and completely metrizable; a standard Borel space is a measurable space isomorphic to the Borel space of a Polish space. The orbit relation is not assumed to be Borel merely because the action is Borel.

\textbf{Question under these conventions.} For every Borel equivalence relation \(E\) on \(\R^\omega\) such that \(E\leq_B E_G^Y\) for some such action, does there exist a sequence \((P_n)_{n\in\omega}\) of nonempty perfect subsets of \(\R\) such that, with
\[
P=\prod_{n\in\omega}P_n,
\]
the restricted relation \(E\mathbin{\upharpoonright}P=E\cap(P\times P)\) is smooth? Explicitly, the requested conclusion is the existence of a Borel map \(f:P\to\R\), for the relative Borel structure on \(P\), satisfying
\[
\forall x,x'\in P,\qquad x\mathrel E x'\ \Longleftrightarrow\ f(x)=f(x').
\]
The supplied statement does not explicitly specify the reducibility notion or the class of group actions meant by ``orbit equivalence.'' The formulation above uses Borel reducibility and Polish group actions as a conventional interpretation requiring confirmation. It interprets ``a product of perfect sets'' as a full product of nonempty perfect coordinate sets; nonemptiness excludes a vacuous conclusion. These conventions are not explicitly defined in the supplied source excerpt.
\subsection{Short English statement}
Under the stated descriptive set theory conventions, must every Borel equivalence relation on \(\R^\omega\) that is Borel reducible to an orbit equivalence relation admit a product of nonempty perfect coordinate sets on which equivalence is classified by equality of a Borel real-valued invariant? The supplied record labels this problem open; that status has not been independently verified.
\subsection{Sources}
\begin{itemize}
\item Supplied local UnsolvedMath problem registry record: permanent ID \(30001678\), title ``Smoothness on Products of Perfect Sets,'' collection ``Oberwolfach Reports Open Problems,'' exact source problem code \texttt{catalogue entry}. The proposed year is recorded as 2011; no individual proposer is identified.
\item Supplied source citation and literature link labeled ``Oberwolfach problem statement'': \emph{Set Theory} (2011), Oberwolfach Reports, pp.\ 85--140, DOI \texttt{10.4171/owr/2011/02}. \url{https://doi.org/10.4171/owr/2011/02}. The background describes these reports as documenting mathematical workshops and their open-problem discussions.
\item Supplied literature link labeled ``canonization on perfect cubes'': \url{https://arxiv.org/abs/1804.05174}. No authors, precise result, or applicability hypotheses are supplied for this reference.
\item Literature assessment date: \texttt{2026-08-21}. This is supplied dated curation, not an independently verified assessment. The record labels the problem open and reports that product-perfect-set canonization is known in special settings while a general smoothness assertion for an arbitrary Borel orbit equivalence relation remains unresolved. The supplied background explicitly requests independent verification of time-sensitive status and cited claims.
\item The supplied statement-status label is \texttt{corrected\_verified}. Its accompanying verification note states only that TeX notation was enclosed in math delimiters without changing mathematical content; it does not establish mathematical correctness or current research status.
\end{itemize}
\end{document}
